% FILE: 02_lineage_and_context.tex
% Project: schemas
% Thesis Role: data-schema-definitions-for-ma-and-receipt-surfaces

\section{Lineage and Context}
\label{sec:schemas-lineage}

\subsection{2016 Language-Model Lesson}
\label{sec:schemas-lineage-2016}

The 2016 language-model experiment established that local text imitation does not conserve
consequence. A model trained to reproduce the surface form of process mining reports can emit
plausible fitness scores, conformance verdicts, and financial projections without those outputs
being grounded in any actual event log, process model, or executing query. The outputs are
structurally indistinguishable from genuine evidence to a reader who does not have access to the
underlying computation.

This lesson is directly relevant to the schemas project. If a receipt is expressed only as a
natural-language or loosely typed JSON document, an adversary can manufacture a receipt that
passes informal inspection without satisfying any of the cryptographic requirements that make
re-execution possible. The schema is the formal response to the 2016 lesson: it moves the
receipt from the domain of surface-form text imitation into the domain of machine-checkable
structural contracts. A receipt that does not carry a 64-hex BLAKE3 hash in \texttt{target\_log\_hash}
will fail schema validation, regardless of how plausible its prose assertions appear.

The schema therefore serves as the boundary between the regime the 2016 experiment identified as
defective---text that imitates evidence---and the regime this research program requires: structured
artifacts that are evidence by construction.

\subsection{Enterprise Process Gap}
\label{sec:schemas-lineage-enterprise-gap}

The enterprise process gap is the observation that enterprise software systems produce operational
data that is not automatically process-intelligence-ready. Event logs must be extracted,
canonicalized, and admitted before they can be used as evidence. Process models must be identified
by content address before they can be cited in a claim. The gap between raw operational data and
board-admissible evidence is where most process intelligence projects fail.

The schemas project addresses the downstream end of this gap. It does not specify how event logs
are extracted or how process models are manufactured---those concerns belong to the \texttt{wasm4pm}
and \texttt{wasm4pm-compat} layers. What it does specify is the shape of the artifact that
\emph{proves} the gap was bridged: a receipt with a BLAKE3 hash of the admitted event log, a
BLAKE3 hash of the discovered process model, a query definition naming the execution engine, and
a signature certifying the result.

Without the schema, the gap-bridging claim cannot be validated. The enterprise produces data;
the process intelligence layer produces a receipt; the schema is the contract that makes the
receipt auditable.

\subsection{Relationship to the Chatman Equation}
\label{sec:schemas-lineage-chatman-equation}

The Chatman Equation $A = \mu(O^*, T, L)$ expresses that an admissible artifact $A$ is the
result of applying the manufacturing function $\mu$ to an objective $O^*$, a type law surface $T$,
and an event log $L$. The schemas project materializes the type law surface $T$ for the receipt
domain: it specifies what type a receipt must have to be admissible.

Concretely, the \texttt{SlideVerificationLedger} schema defines the type of the receipt artifact.
A receipt of type \texttt{SlideVerificationLedger} is one that carries a \texttt{slide\_id} of
type UUIDv4, a \texttt{target\_log\_hash} of type 64-hex BLAKE3, a \texttt{process\_model\_hash}
of type 64-hex BLAKE3, a \texttt{query\_definition} with an \texttt{engine} field constrained to
the \texttt{wasm4pm}~$|$~\texttt{pm4py} enum, a \texttt{verification\_results} object with a
status in \{\texttt{verified}, \texttt{failed}, \texttt{unverified}\}, and a
\texttt{validator\_signature} of type 128-hex Ed25519.

The manufacturing function $\mu$ must produce artifacts of this type. Artifacts that do not
conform are not of type $A$---they are rejected by the schema's \texttt{additionalProperties: false}
constraint, which refuses any receipt that carries fields outside the defined property set.

\subsection{Relationship to Receipts and Replay}
\label{sec:schemas-lineage-receipts-replay}

The receipt-and-replay architecture requires that every process intelligence claim be backed by a
receipt that enables independent replay of the underlying analysis. Replay requires three
identifiers: the event log (by content address), the process model (by content address), and the
query (by URI and engine). The \texttt{SlideVerificationLedger} schema encodes exactly these three
identifiers as required fields, plus the results of the original execution and a signature that
binds the whole payload together.

The \texttt{OTelTraceBlake3IntegrationSchema} extends the receipt architecture to the execution
trace level. Where the ledger schema tracks the content of the \emph{inputs} to the analysis
(event log, process model), the OTel schema tracks the content of the \emph{execution} itself:
each span in the trace carries a \texttt{blake3\_receipt} computed as
$\operatorname{BLAKE3}(\text{PriorReceipt} \| \text{SpanData})$, forming a chain that makes the
execution unforgeable. A forged trace would require inverting the BLAKE3 hash function at every
span boundary---a computationally infeasible operation.

Together, the two schemas cover both the content of the analysis and the execution path that
produced it. This two-layer coverage is what distinguishes receipt-backed process intelligence
from informal reporting: the claim is not only asserted, it is traced back to an execution that
can be replayed and verified.
